\documentclass[12pt]{article}
\usepackage{amsmath, amssymb}
\usepackage[margin=1in]{geometry}
\setlength{\parskip}{6pt}
\title{QUBO Formulation: QUBO Feature Selection}
\date{}
\begin{document}
\maketitle

\subsection*{QUBO Feature Selection}

\paragraph{Problem Statement.}
Let $N$ be the number of candidate features, and let $\mathcal{F} = \{0, 1, \dots, N-1\}$ denote the set of features. Each feature $i \in \mathcal{F}$ has an associated relevance score $s_i \in \mathbb{R}$. For every pair of distinct features $i, j \in \mathcal{F}$, there is a redundancy penalty $p_{ij} \in \mathbb{R}$, where $p_{ij} = p_{ji}$. The goal is to select a subset of features of cardinality exactly $K$ such that the total relevance score of the selected features minus the sum of redundancy penalties between all pairs of selected features is maximized.

\paragraph{Decision Variables.}
For each feature $i \in \mathcal{F}$, we introduce a binary decision variable $x_i \in \{0, 1\}$. The variable $x_i$ takes the value $1$ if feature $i$ is selected, and $0$ otherwise.

\paragraph{QUBO Derivation.}
\textbf{Step 1: Objective Function.}
The original maximization objective is
\begin{equation}
    \max_{x} \sum_{i \in \mathcal{F}} s_i x_i - \sum_{\{i,j\} \subseteq \mathcal{F}, i < j} p_{ij} x_i x_j.
\end{equation}
Following QUBO convention, we convert this to a minimization problem by negating the objective:
\begin{equation}
    \min_{x} -\sum_{i \in \mathcal{F}} s_i x_i + \sum_{\{i,j\} \subseteq \mathcal{F}, i < j} p_{ij} x_i x_j.
\end{equation}

\textbf{Step 2: Cardinality Constraint.}
The requirement to select exactly $K$ features is expressed as the equality constraint
\begin{equation}
    \sum_{i \in \mathcal{F}} x_i = K.
\end{equation}

\textbf{Step 3: Penalty Term Expansion.}
We enforce the constraint via a quadratic penalty term with weight $\lambda > 0$:
\begin{equation}
    P(x) = \lambda \left( \sum_{i \in \mathcal{F}} x_i - K \right)^2.
\end{equation}
Expanding the square yields
\begin{equation}
    P(x) = \lambda \left( \left(\sum_{i \in \mathcal{F}} x_i\right)^2 - 2K \sum_{i \in \mathcal{F}} x_i + K^2 \right).
\end{equation}
Using the identity $\left(\sum_{i \in \mathcal{F}} x_i\right)^2 = \sum_{i \in \mathcal{F}} x_i^2 + 2 \sum_{\{i,j\} \subseteq \mathcal{F}, i < j} x_i x_j$ and the idempotent property $x_i^2 = x_i$, we obtain
\begin{equation}
    \left(\sum_{i \in \mathcal{F}} x_i\right)^2 = \sum_{i \in \mathcal{F}} x_i + 2 \sum_{\{i,j\} \subseteq \mathcal{F}, i < j} x_i x_j.
\end{equation}
Substituting this back into the penalty expression and collecting terms gives
\begin{equation}
    P(x) = \lambda \sum_{i \in \mathcal{F}} (1 - 2K) x_i + 2\lambda \sum_{\{i,j\} \subseteq \mathcal{F}, i < j} x_i x_j + \lambda K^2.
\end{equation}

\textbf{Step 4: Combined Hamiltonian.}
Adding the negated objective and the penalty term, the QUBO Hamiltonian $H(x)$ is
\begin{equation}
    H(x) = \sum_{i \in \mathcal{F}} \left( \lambda(1 - 2K) - s_i \right) x_i + \sum_{\{i,j\} \subseteq \mathcal{F}, i < j} \left( p_{ij} + 2\lambda \right) x_i x_j + \lambda K^2.
\end{equation}

\textbf{Step 5: Q Matrix and Offset.}
The Hamiltonian is of the form $H(x) = x^T Q x + c$. The entries of the upper-triangular Q matrix and the offset $c$ are given by
\begin{align}
    Q_{i,i} &= \lambda(1 - 2K) - s_i, \quad \forall i \in \mathcal{F}, \\
    Q_{i,j} &= p_{ij} + 2\lambda, \quad \forall i, j \in \mathcal{F}, i < j, \\
    c &= \lambda K^2.
\end{align}

\paragraph{Penalty Weight Justification.}
Let $S_{\max} = \max_{i \in \mathcal{F}} s_i$. The penalty weight $\lambda$ must satisfy $\lambda > S_{\max}$. This condition ensures that the energy cost incurred by violating the cardinality constraint exceeds the maximum possible gain in relevance score achievable by any single feature. Specifically, the quadratic penalty grows sufficiently fast to dominate the linear relevance terms, guaranteeing that the global minimum of $H(x)$ corresponds to a feasible solution satisfying the cardinality constraint.

\paragraph{Correctness Argument.}
Minimizing $H(x)$ over $x \in \{0, 1\}^N$ recovers the optimal feature selection. For any feasible solution where $\sum_{i \in \mathcal{F}} x_i = K$, the penalty term vanishes, and $H(x)$ equals the negated objective function. For any infeasible solution, the penalty term is strictly positive. Due to the choice of $\lambda$, the energy of any infeasible solution is strictly greater than the energy of the best feasible solution. Consequently, the global minimum of $H(x)$ must be feasible and minimizes the original objective, yielding the optimal set of $K$ features.

\paragraph{Illustrative Example.}
Consider the concrete instance with $N=3$ features, so $\mathcal{F} = \{0, 1, 2\}$. The general formulas instantiate as follows. The Hamiltonian is
\begin{align}
    H(x) &= \sum_{i=0}^{2} \left( \lambda(1 - 2K) - s_i \right) x_i + \sum_{0 \le i < j \le 2} \left( p_{ij} + 2\lambda \right) x_i x_j + \lambda K^2 \\
    &= \left( \lambda(1 - 2K) - s_0 \right) x_0 + \left( \lambda(1 - 2K) - s_1 \right) x_1 + \left( \lambda(1 - 2K) - s_2 \right) x_2 \nonumber \\
    &\quad + \left( p_{01} + 2\lambda \right) x_0 x_1 + \left( p_{02} + 2\lambda \right) x_0 x_2 + \left( p_{12} + 2\lambda \right) x_1 x_2 + \lambda K^2.
\end{align}
The Q matrix entries are
\begin{align}
    Q_{0,0} &= \lambda(1 - 2K) - s_0, & Q_{0,1} &= p_{01} + 2\lambda, \\
    Q_{1,1} &= \lambda(1 - 2K) - s_1, & Q_{0,2} &= p_{02} + 2\lambda, \\
    Q_{2,2} &= \lambda(1 - 2K) - s_2, & Q_{1,2} &= p_{12} + 2\lambda,
\end{align}
with offset $c = \lambda K^2$. The optimal selection is the binary vector $x^*$ that minimizes $H(x)$, corresponding to the subset of $K$ features with the highest net score.
\end{document}